\subsubsection{Index Matching}
The predicates in the queries can contain comparison operators, as well as \texttt{AND} and \texttt{OR}, allowing conjuncts and disjuncts.

Index Matching is the process of choosing the proper index for the task.
For conjuncts only, we can use:
\begin{itemize}
    \item \bi{Hash Indexes}. Constraints: Only equality operations and predicate has to contain all index attributes
    \item \bi{B+-Tree}. Constraints: The attributes in the predicates need to be prefixes of the search key
\end{itemize}

Here, the indexes need to match \bi{every} predicate in the disjunction.
For instance, if we have a hash index on $A$ and a B+ tree on $B$, and the query is \verb|A = 7 OR B > 5|.
In that case, we can either file scan or use both indexes (we fetch the RIDs and take the union)

However, index matching is a hard problem.
